\chapter{回到三平均：几种终止方法放在一起看}\label{ch:three}
\section{最小的奇素数，不能简单代入所有大p公式}
三平均的最终阈值是七。七元已被奇数余量构造覆盖，九元可排成 $3\times3$，先平均三行，再平均三列。因此真正还需处理的是八元、十元和所有十一元以上。

它们展示三种不同的终止方法：十一元以上保持整数平方和下降；八元让一个本原参数的绝对值下降；十元让一个由局部类型决定的整数二次量下降。有限步证明并不一定要所有分支用同一个下降量。

\section{十一元以上：三个余数类使分类更直接}
模三只有三个类。有四个不同三重值，就必有两个同余；有三个三重值，就看它们是否各占一类。很多一般 $p$ 的长选择，在这里变成 $(a,a,x)$ 或 $(a,b,x)$。

最细的地方是危险素数可能为二。如果选择 $(a,x,x)$，只靠“还留着一份 $a$”不能从 $2x\equiv2a\pmod2$ 推出 $x\equiv a$。证明会额外使用碰撞已经给出的 $x=3b-2a$，并留下 $a,b$ 的副本。这个\emph{已知的整数仿射关系}修复了模二失效。

全部分类的剩余容量来自 $n\ge d_3(n)+10$。所以一旦需要保护最多 $d_3(n)$ 个位置，仍有足够轻值可选。
\InsertProof{lem:triple-large}

\section{八元：从加减倍数回到欧几里得算法}
八元先进入 $(u^4,v^3,-4u-3v)$。取 $(x,y)=(u,u+v)$ 后，合法性是 $y$ 奇。目标是把正整数 $|y|$ 降到一。

两条真实返回提供 $x\mapsto x\pm2y$。这让我们能在一个模 $2y$ 的剩余类里选一个离 $2y$ 最近的 $x_0$。第三条返回把第二坐标改成 $2(x_0-2y)$；两个输出都有公因子二，约去以后，新的 $|y|$ 比原来小。这和欧几里得算法“减去合适倍数后取小余数”的逻辑一样。

例如分析列 $(x,y)=(1,3)$。加一次 $2y$，选 $x_0=7$，此时 $|7-6|=1<3$。施加第三返回得到 $(30,2)$，取本原代表是 $(15,1)$。一轮便把 $|y|$ 从三降到一。选择平移的准则很具体：让 $x_0$ 尽量接近 $2y$，以便下一步的第二坐标变小。
\InsertProof{lem:eight}

\section{十元：有限表为什么可以覆盖无限有理输入}
十元的危险素数只有二和五。先用安全整数下降取得六重值，再经一个四步桥进入 $(u^6,v^3,-6u-3v)$。

取方向坐标 $(X,Y)=(u,3u+2v)$ 后再取本原代表。不同奇偶类型对应不同的尺度 $h=1$ 或 $2$，势能写成
\[h^2(5X^2+Y^2).
\]
不要把 $h$ 忽略：同一形式的有理方向，在清分母与取本原代表时可能需要不同的整数尺度，比较下降必须用同一真实规则。
\InsertProof{lem:ten}

\section{具体看一个二次下降条件}
设输入属于 $h=1$ 的一类，采用
\[M_4=\begin{pmatrix}1&1\\5&-3\end{pmatrix},
\]
该类的输出内容为四，输出尺度仍为一。乘掉内容分母后，输入减输出的量为
\[
16(5X^2+Y^2)-[5(X+Y)^2+(5X-3Y)^2]
=50X^2+20XY+2Y^2.
\]
它等于 $2(5X+Y)^2$。所以除了直线 $Y=-5X$，这张矩阵都严格下降。在那条直线上不能就停止，还要检查同一局部类型的其他矩阵；完整表做的正是多个下降区域的联合覆盖。

原问题的有理数可以任意大，但决定内容和模四类型只需要有限余数。固定一种类型后，是否下降变成实变量 $t=Y/X$ 上的二次不等式。二次函数只在至多两个根处变号，因此可精确检查全部实区间，进而覆盖其中所有有理方向。
\InsertProof{prop:ten-certificate}
\InsertTables{16_three.tex}
表里的 $\infty$ 表示 $X=0$，这正是含零终点。有限未覆盖点 $0$ 表示 $Y=0$，违反 $5\nmid Y$。核验应同时检查“没有下降”与“这个方向是否真的属于该局部类型”；不能仅凭一个实图上看起来被盖住了就结束。
